\documentclass[openany,fontset=fandol]{ctexbook}
\usepackage{amsmath,amssymb}
\usepackage[hidelinks]{hyperref}
\newcommand{\note}[1]{\footnote{译注：#1}}
\title{代数几何基础第五卷预备笔记\\第12节：记号与两种纤维结构}
\author{据用户上传英译编辑稿译出}
\date{2026年11月16日}
\begin{document}
\maketitle
\chapter*{范围与校订状态}
本批连续译出底本第44页“I use [or profit from] this opportunity for a notational self-criticism”起的四段，止于“$X^{(m)}$ over $\operatorname{Gr}^{m}$”这一括注末。下一处是同页“The Proposition 2.1”段。

本批沿用上一批的 $\operatorname{rank}E=r$、上标 $m$ 表示余维的约定。底本的纤维维数公式与此不相容，正文按核模的秩校正，并在脚注保留原式与依据。记号试拟不当作已证明同构。此前缺文和待核事项未解决；本批未合入原合订本。
\chapter*{\S12 将上述结果推广到线性截面（续）}
\section*{关于记号的自我批评}
我借此机会对记号作一番自我批评，这段话也可以放在第1节。问题在于：对于与 $X$ 相对应的对象，任意选取字母 $Y$ 来表示，也就是把字母 $Y$ 与 $X$ 配在一起。这样，一旦 $X$ 改成 $Z$，就不大清楚应该改用什么字母了。这种不便已经使我采用了一些不一致的记号。

也许，这里引入整数 $m$ 的较一般情形，提示了一个合理的办法：用 $X^{(m)}$ 代替 $Y^{(m)}$，于是第1节中的 $Y$ 就改写为 $X^{(1)}$。如此，迭代截取的记号或许能与余维相加的记法大体协调。\note{底本在此试拟“we might approximately have $X^{(m)(m_1)}=X^{m+m_1}$”。这是关于记号的愿望，不能直接当作两个普遍族的典范同构：连续选取两个线性截面通常保留先后选择或旗标的数据，而一次选取总余维截面只保留最终子空间；二者参数空间一般不同。例如在射影平面中先选直线再选其上一点，参数为点线旗标；直接选余维2子空间只参数化点，忘掉直线的映射具有正维纤维。因此正文只陈述记号意图，原等式保留于此，不把它升格为定理。}
我将在下文尝试采用这种记号。显然，这个上标本身也可以受到批评，因为代数几何中常用上标表示所讨论簇的维数。不过，我们绝不采用那一种约定，因此我认为在这一点上可以自由选择。

\section*{作为格拉斯曼预概形的结构}
在前述 $X^{(m)}$ 的构造中，有典范同构
\[
X^{(m)}\simeq\operatorname{Grass}_{m}(F_X^\vee),
\]
其中，若 $f:X\to P$ 是给定态射，则
\[
F_X=\ker\bigl(E_X\twoheadrightarrow\mathcal O_X(1)\bigr)
\simeq f^*\bigl(\Omega^1_{P/S}(1)\bigr).
\]
\note{底本写作“$F_X\simeq F^*(\Omega^1_{P/S})(1)$”，把态射与核模的字母混用。这里明确以 $f$ 表示态射，以 $F_X$ 表示核模。射影丛的Euler正合列拉回后仍正合，因为右端可逆且正合列局部分裂。因此 $F_X$ 局部自由，秩为 $r-1$。上述同构也可直接按函子核验：关联条件是泛子模 $G^\vee\subset E_X$ 的像落入 $F_X$；它在 $F_X$ 中仍是局部直和因子，因为 $F_X/G^\vee$ 是满射 $E_X/G^\vee\to\mathcal O_X(1)$ 的局部自由核。对偶后即得到 $F_X^\vee$ 的秩 $m$ 局部自由商，反向构造相同，并与基变换相容。这里的说明是译者补充，不是底本另列的证明。}
特别地，当 $0\leq m\leq r-1$ 时，$X^{(m)}$ 在 $X$ 上光滑，且其几何纤维为整概形。事实上，纤维是维数
\[
m(r-1-m)
\]
的有理簇。\note{底本原式为 $m(r-m)$。但本节上一页已重新规定 $\operatorname{rank}E=r$，所以核模 $F_X$ 的秩是 $r-1$；其秩 $m$ 商的格拉斯曼簇维数为 $m((r-1)-m)$。例如 $r=2,m=1$ 时，射影直线上包含给定点的超平面只有该点，纤维维数为0，不是原式给出的1。若取 $m=r$，该格拉斯曼概形为空，不能无条件声称纤维几何整；故正文明确范围。端点 $m=0,r-1$ 的纤维为一点。中间范围内，矩阵坐标图给出光滑性及有理开图，几何整性是通常格拉斯曼簇的性质。}
当然，验证可以归约到 $X=P$ 的情形；因此，和前面的讨论一样，这属于关于格拉斯曼预概形的一般事实。我相信你会在另一个段落中将这些一般事实充分展开。

\section*{在参数空间上的射影丛结构}
现在得到了与第1节中的图式完全类似的情形。还有一点此前忘了提：作为 $\operatorname{Gr}^{m}$ 上的预概形，$H^{(m)}$ 典范地同构于
\[
\mathbf P\bigl(E_{\operatorname{Gr}^{m}}/G^\vee\bigr).
\]
因而，在 $0\leq m\leq r-1$ 的范围内，它是一个很好的射影丛。\note{这里 $G$ 是上一批中的 $E^\vee_{\operatorname{Gr}^{m}}$ 的泛秩 $m$ 商，所以 $G^\vee$ 是 $E_{\operatorname{Gr}^{m}}$ 的局部直和因子，余商的秩为 $r-m$，射影丛相对维数为 $r-m-1$。秩为零的端点 $m=r$ 给出空预概形，故不将其称为非空射影纤维。与前段不同，这里讨论的是 $H^{(m)}\to\operatorname{Gr}^{m}$，而不是 $X^{(m)}\to X$。}
当然，一般不能由此对 $\operatorname{Gr}^{m}$ 上的 $X^{(m)}$ 得出同样的结论。\note{底本在括号中明确保留这一限制。$X^{(m)}=X\times_P H^{(m)}$ 的参数纤维是 $X$ 的线性截面，可能奇异、可约或为空；本批不把前段关于另一投影的光滑性误用于这里。}
\end{document}
